% Requires: \usepackage{booktabs} \usepackage{array} \usepackage{multirow}
% (tabularx no longer needed -- explicit p{} widths give more control over the
% Label/Definition/Example proportions than X columns splitting evenly would.
% multirow is needed for the merged Stage 1+2 table's Stage column.)
%
% Each table below is a `table*` (spans both columns of the two-column layout,
% giving the full \textwidth to work with instead of one narrow column) and
% uses fixed p{} widths sized so Definition gets most of the room and Example
% less. `table*` floats can only land at the top/bottom of a page (or a float
% page) in two-column mode -- they will drift slightly from sitting literally
% under each itemize bullet, which is normal/expected for wide floats in a
% two-column document.

\subsection{Theoretical Framework and Codebook}
To identify racialised discourse, which we anticipate is a rare occurrence in the reviews, we created a hierarchical codebook to capture it.
First, we are only interested in racialisation that is tied to social encounters at the pool. We define a social encounter review as one that references other people --- users or staff --- their presence, behaviour, appearance, language, or the atmosphere they create at the pool. The label is social encounter, and it is binary.

Subsequently, we classify reviews that fall in the social encounter class as racialised or not. A review is racialised if --- reading it as a whole --- it attributes a behaviour, characteristic, attitude, presence, or quality of a group defined by race, ethnicity, nationality, culture, or religion.

% ------------------------------------------------------------------
% Stage 1 + 2 -- Social Encounter and Racialised (merged, single-column)
% ------------------------------------------------------------------
% Single-column `table` (not `table*`) so this one sits in \linewidth rather
% than spanning both columns -- footnotesize (one step below \small) and a
% combined "Stage: Label" column (no multirow) to keep it fittable; Definition
% still wraps to several lines per row given how narrow a single column is.
\begin{table}[htbp]
\centering
\footnotesize
\setlength{\tabcolsep}{3pt}
\caption{Stage 1 -- Social Encounter (\texttt{social}) and Stage 2 -- Racialised (\texttt{racialised}, applied to all Stage-1 Yes reviews).}
\label{tab:codebook-social-racialised}
\begin{tabular}{@{} l p{0.35\linewidth} p{0.24\linewidth} @{}}
\toprule
Label & Definition & Example \\
\midrule
Social: Yes & References other people -- users or staff -- their presence, behaviour, appearance, language, or the atmosphere they create. & ``The lifeguard was rude to me'' \\
\addlinespace
Social: No & Only the physical environment, facilities, pricing, or the reviewer's solitary experience; no reference to other people. & ``The pool water was very cold'' \\
\midrule
Racialised: Yes & Attributes anything to a group defined by race, ethnicity, nationality, culture, or religion -- explicitly, implicitly, or coded; any tone. & ``Too many Arabs -- doesn't feel like home'' \\
\addlinespace
Racialised: No & No racial/ethnic/national/religious anchor -- classist, gendered, generic-conduct, or tourist complaints that would read identically about any crowd. & ``Cheap entrance attracts a certain clientele'' \\
\bottomrule
\end{tabular}
\end{table}

Now, in the set of racialised reviews, we study racialisation across four non-hierarchical dimensions.
\begin{itemize}
    \item \textbf{Theme} -- What object or domain is being racialised? Theme follows the object of racialisation --- what/who is being racialised. We define nine themes -- Staff, Crowd, Belonging, Safety, Hygiene, Etiquette, Access, Language, Clothing, Security and Other to capture any other theme. A review could have up to three themes, if two or more distinct objects are genuinely and independently racialised.
    \item \textbf{Context} -- Was the racialisation expressed as a specific incident, a generalisation, or both?
    \item \textbf{Approach} -- How is the racialised group encoded in the review? This dimension classifies the encoding strategy used for each distinct racialised group.
    \item \textbf{Connotation} -- What is the tone of the racialised content? Note: capture the tone of the racialised content only --- not the overall review sentiment.
\end{itemize}

% ------------------------------------------------------------------
% Stage 3 -- Theme, Context, Approach, Connotation (merged)
% ------------------------------------------------------------------
% One table for all four dimensions (Dimension column via \multirow), since
% each dimension is already defined above in prose -- the caption here just
% names the table, it doesn't redefine each dimension again.
\begin{table*}[htbp]
\centering
\small
\caption{Stage 3 -- label definitions and a qualifying example for each of the four content dimensions (\texttt{theme}, \texttt{context}, \texttt{approach}, \texttt{connotation}), applied only to reviews already classified as racialised.}
\label{tab:codebook-dimensions}
\begin{tabular}{@{} l l p{0.46\textwidth} p{0.24\textwidth} @{}}
\toprule
Dimension & Label & Definition & Example \\
\midrule
\multirow[t]{11}{*}{Theme} & Staff & Pool staff are racialised as individuals or as a group -- their composition, features, conduct, attitude, behaviour, or language. & ``The Slavic-sounding lifeguard is disrespectful'' \\
 & Crowd & Pool users are racialised -- who is present, absent, or how the pool population is characterised in racial/ethnic/national terms. & ``Beautiful pool, but there are too many Arabs'' \\
 & Belonging & Racialisation of cultural fit, acceptance, or entitlement to be at the pool -- felt by the reviewer or projected onto others. & ``As a Muslim female, I did not feel accepted here'' \\
 & Language & Racialisation of what language is spoken (or not), by whom -- staff, clientele, or signage. & ``All you hear at this pool is fluent Turkish'' \\
 & Clothing & Swimwear or attire is racialised; includes unequal dress-code enforcement along racial, cultural, or religious lines. & ``I was not allowed to wear my burka'' \\
 & Hygiene & Cleanliness of people, bodies, or sanitary practices is racialised. & ``Dirty Turkish people everywhere'' \\
 & Etiquette & Behaviour, attitude, or adherence to pool norms is racialised. & ``They do not know German sauna culture'' \\
 & Security & Security procedures or enforcement are racialised -- who is checked, how, and whether applied equally. & ``The guard searched us multiple times, but not the Germans ahead of us'' \\
 & Safety & Personal safety or threat perception is racialised -- feeling unsafe or attributing danger/criminal behaviour to a racialised group. & ``An Arab boy stole my daughter's toy'' \\
 & Access & Entry or access is racialised -- denied, restricted, delayed, or informally excluded on racial/ethnic/cultural/religious grounds. & ``We were denied access because of our skin colour'' \\
 & Other & Racialisation is present and attributable but fits none of the above categories. & -- \\
\midrule
\multirow[t]{2}{*}{Context} & Incident & Racialisation is attached to a specific event, interaction, or occurrence -- something happened, with a describable moment or action. & ``The receptionist was racist when she found out I am Ukrainian'' \\
 & Generalised & Racialisation is a generalisation or standing judgement about a person or group, not tied to a specific event; includes atmospheric observations about the pool in general. & ``Turkish people don't follow the rules'' \\
\midrule
\multirow[t]{3}{*}{Approach} & Explicit & Directly identifiable by a general reader without inference or specialist knowledge. \emph{Would a general reader immediately understand which group is referenced?} & ``There are too many Turkish people at the pool'' \\
 & Coded & Identifiable only to a reader with specific in-group, cultural, or political knowledge; deliberate, deniable to outsiders. \emph{Would an in-group reader know who is meant, while a moderator might not?} & ``There aren't as many Ayse's and Fatma's here'' \\
 & Implicit & No specific group identifiable from the text, even by an in-group reader; vague framing -- atmosphere, nostalgia, unnamed attribution. \emph{Could any reader identify the group from the text alone?} If no, Implicit. & ``There is no room for us here anymore'' \\
\midrule
\multirow[t]{3}{*}{Connotation} & Positive & Racialisation is favourable -- positive stereotypes, flattering attributions, or appreciation linked to racial/ethnic/cultural identity. & ``The German lifeguard is so patient with the kids'' \\
 & Negative & Racialisation is unfavourable -- an us-vs-them dynamic, harm/threat attributed to a group, or discrimination framed in racial terms. & ``Dirty Turkish people everywhere -- no hygiene culture'' \\
 & Neutral & Racialisation is present but tone is indeterminate -- a factual/demographic observation with no qualifying adjective conveying approval or disapproval. & ``The lifeguard speaks English and German'' \\
\bottomrule
\end{tabular}
\end{table*}
